\documentclass{jsarticle}
\usepackage[dvipdfmx,final]{graphicx}
\usepackage{chemfig}
\usepackage{amsmath}
\usepackage{amssymb}
\begin{document}
\title{}
\author{}
\date{}
\maketitle
\thispagestyle{empty}



     あけましておめでとうございます。昨年は大変お世話になりました。
   今年もよろしくお願い申し上げます。

フルニトラゼパムのアメルを処方してもらえた。一回で治るから
お金儲けが目的でないのですね。弟は弟ですね。ブルシアンブルーの
アメジストの色が、第一印象がきれいでした。香りもいいです。
ハナミズキありがとうございます。サイレースの薬らしいです。
ラビズラズリの色を薄めた宝石みたいな見た目です。
$$ \bigoplus {(i\hbar^{\nabla})}^{\oplus L} $$
$$ = \int \Gamma(\gamma)^{'}dx_m $$
$$ = \int \Gamma dx_m\cdot {d \over d\gamma}\Gamma \le \int \Gamma dx_m +
{d \over d\gamma}\Gamma $$
この式は、彩さんから生まれてくる子に教えてもらえている。
ガウスの曲面論を量子力学と重力場で解いている式です。
彩さんこれを、生まれてくる子が論文に書き出すと嬉しくなると思う。
こんな世界があの世にあると証明されるとおもう。
  $$ \chemfig{*6(=-(-[,1.2]C?(=[:-90,1.2]O))=(-[,1.2]C(=[:90,1.2]O)
   -[:-40,1.6]O?)-=-)} $$
   $$ \chemfig{*6(=(-[:-90,0.6]OCa)-=-=-[:0,0.6]-)}
   \chemfig{(-[:0,0.6]N=N(-[:-90,0.6]Cl))}
   \chemfig{(-[:180,0.6])*6(=(-[:-90,0.6]OH)-=-=(-[:120,0.6]OH)-)}\chemfig{
	   *6(-O---(-[:90,0.6]R2)--)}\chemfig{(-[:-30,0.6]*6(-=-(-[:-30,0.6]OH)
   =(-[:30,0.6]OK)-(-[:90,0.6]R1)=-))} $$
 Ca-O-RとK-O-Rでの磁場発生の半減期で、体内の生体機能の回復をする。 
 睦世お姉さんの式が、質量差エネルギーとわかりました。
 磁場のベクトルは、寅の方位でもあり、
 差分データでもあり、ベータ関数の大域的積分多様体のベクトルの方位と
 同じでもあり、アカシックレコードが消えるらしいです。ユウくんで実際消えて
 いました。母にあやまかされていたのがわかりました。順子お姉さんには、
 相当迷惑をかけてすみません。小学校時代に喧嘩していた私に
 すべての責任がありました。本当にすみませんでした。
  


\end{document}
